\section{相变、潜热与内能变化}
\subsection{由相态确定热量计算的阶段}
同一种物质可处于固态、液态或气态。在给定温度、压强下，哪种相稳定取决于物质的性质和外部条件。液体与蒸气可以在一定条件下共存，固体与液体也可共存；“同一物质”不等于任何时候都只有一个相。

常见的一阶相变如熔化、汽化，在给定压强下可以近似恒温进行，同时伴随潜热和体积变化。这里的恒温并不是所有相变的无条件特征：外压改变、混合物成分改变、成核迟滞或接近临界点时，都需要重新讨论。

在远离相变的单相区，热量通常写成 $Q=mc\Delta T$ 或 $nC\Delta T$；在恒温相变区，温度不变却仍需交换能量，必须另计潜热。把全过程只写成一个 $mc\Delta T$ 会漏掉相变能量。


每千克潜热记为 $L$，每摩尔潜热记为 $L_m$，两者满足 $L_m=\mu L$。熔化、汽化取吸热为正，凝固、凝结取放热为负。使用数值时应先核对潜热是按质量还是按物质的量给出，不能把kg与mol混用。

若物质经历升温、相变、再升温，应按实际阶段相加。例如在恒定外压下，液体先从 $T_i$ 升至沸点 $T_b$，再全部汽化，所需热量为 $mc_l(T_b-T_i)+mL_v$；若还要把蒸气升温，须再加蒸气的显热项。

\fig{fig:phase-heating}{恒定外压下加热跨越汽化区的分段关系}{\begin{tikzpicture}[x=1cm,y=1cm]
\draw[axis] (0,0)--(5.4,0) node[right]{$Q$};\draw[axis] (0,0)--(0,3.5) node[above]{$T$};\node[below left] at (0,0){$O$};\draw[ray] (.5,.6)--(2,1.7);\draw[ray] (2,1.7)--(4,1.7);\draw[ray] (4,1.7)--(5,2.8);\draw[guide] (0,1.7)--(2,1.7);\node[left] at (0,1.7){$T_b$};\node[above] at (1.1,1.2){液体升温};\node[above] at (3,1.7){气液共存};\node[right] at (4.5,2.9){蒸气升温};\node[below] at (3,1.4){$\Delta Q=mL_v$};
\end{tikzpicture}}

图\ref{fig:phase-heating}横轴是累计输入热量，平台的长度由相变量和潜热决定。若输入热量只够汽化一部分，末态停在平台内，应由剩余能量求汽化质量，而不是直接进入蒸气升温段。

平台上的温度不变，并不表示输入能量没有作用，而是能量改变了两相的比例及分子间相互作用能，同时可能推动外界。图中任一点都应附带“哪一相存在、各有多少”的信息，单写温度不能完整确定处于平台内的状态。因此相变题除了温度方程，还需要一条物质分配关系。例如水蒸气凝结后，总水量仍守恒，气相量却减少，不能继续把最初的蒸气摩尔数固定在状态方程中。

\subsection{潜热为何不等于内能差}
恒压可逆汽化时，液体吸收的潜热既用于改变内能，也用于推动外界。对1mol物质，
\[L_m=\Delta u+p(\bar V_v-\bar V_l)=\Delta h.
\]
其中 $\bar V_v,\bar V_l$ 是气、液两相的摩尔体积，$h=u+p\bar V$ 为摩尔焓。若蒸气可按理想气体处理、液体体积相对很小，$p\bar V_v=RT$，于是

\[\boxed{\Delta u_{\text{汽化}}\simeq L_m-RT}.
\]
\fig{fig:h4:latentwork}{恒温恒压汽化时活塞向外移动；潜热同时供给内能增加和膨胀功}{\begin{tikzpicture}

\foreach \j/\hh/\name in {0/.35/液体,1/3/蒸气}{\begin{scope}[xshift={4.8*\j*1cm}]\draw[world](0,3.5)--(0,0)--(1.8,0)--(1.8,3.5);\fill[blue!12](0,0)rectangle(1.8,\hh);\draw[very thick](0,\hh)--(1.8,\hh);\node at(.9,-.6){1mol\name};\node at(.9,4.1){外压$p$};\end{scope}}
\draw[guide](4.6,.35)--(7.5,.35);\draw[dim](7.5,.35)--(7.5,3)node[midway,right]{$\Delta x$};\draw[flow](2.4,1.8)--(4,1.8)node[midway,above]{汽化};\node at(3.4,-1.4){$L_m=\Delta u+p\Delta\bar V$};

\end{tikzpicture}}
图\ref{fig:h4:latentwork}中，活塞承受同一外压，但气、液的摩尔体积不同。向外推动的功为$p\Delta\bar V$，它已包含在恒压吸热$L_m$中，所以不能把$L_m$全算作内能增加。凝结沿反向发生时，外界压入的功为正，放出的热量除内能减少外，还包含这部分输入功。



凝结时热量和体积功都取相反符号，内能变化也反号。例\ref{t4-1}中放热量的绝对值大于内能减少量，是因为外界在凝结压缩时对系统做了正功。

这一区别也说明恒压潜热不能不加判断地用于刚性容器。刚性容器边界功为零，但相变后的温度、压强和气液比例可能与恒压过程不同；只有先确定实际始末态，才可用 $Q=\Delta U$ 求热量。


潜热的数值还必须与质量单位对应。比潜热按每千克定义，摩尔潜热按每摩尔定义，乘入的相变量也相应不同。只有实际发生相变的那部分质量乘潜热；已经成为过热蒸气的部分继续升温，还需另计显热。将全部初始质量都乘一次汽化潜热，可能掩盖末态仍有液体或已经完全汽化的区别。


焓差是状态差，恒压过程中的吸热等于焓差还需要相应的功条件。把这一性质与装置边界结合起来，才能把恒压测得的相变数据用于其他路径：先将潜热换成所需的内能差，再按实际边界计算功和热量。

\subsection{冰水混合中的相态判断}
设质量 $m_i$、初温 $T_i<T_m$ 的冰与质量 $m_w$、初温 $T_w>T_m$ 的水混合，熔点为 $T_m$。外压恒定，忽略热损失、量热器热容和其他机械能变化，冰、水比热分别为常量 $c_i,c_w$，熔化潜热为 $L_f$。先把两部分都与熔点状态比较，定义

\[E=m_wc_w(T_w-T_m)-m_ic_i(T_m-T_i).
\]
$E$ 表示热水冷却到熔点可放出的热量，扣除冰升到熔点所需热量后的差额。这个差额决定是否还需凝固一部分水，或能够熔化一部分冰，不能一开始就假设冰已全部熔化。

\fig{fig:ice-water}{冰水混合的相态由能量差额所在区间确定}{\begin{tikzpicture}[x=1cm,y=1cm]
\draw[axis] (0,0)--(10,0) node[right]{$E$};\draw (2,-.15)--(2,.15) (5,-.15)--(5,.15) (8,-.15)--(8,.15);\node[below] at (2,-.15){$-m_wL_f$};\node[below] at (5,-.15){$0$};\node[below] at (8,-.15){$m_iL_f$};\node[align=center] at (.8,1){全部为冰\\$T_f<T_m$};\node[align=center] at (3.5,1){部分水凝固\\$T_f=T_m$};\node[align=center] at (6.5,1){部分冰熔化\\$T_f=T_m$};\node[align=center] at (9.3,1){全部为水\\$T_f>T_m$};
\end{tikzpicture}}

图\ref{fig:ice-water}中的四个区间对应四种结果。若 $-m_wL_f\le E<0$，水需凝固 $-E/L_f$ 的质量，末温为 $T_m$；若 $0\le E\le m_iL_f$，冰熔化 $E/L_f$，末温也为 $T_m$。两种情况都要保留未发生相变的剩余部分。

若 $E<-m_wL_f$，全部水凝固仍不足以把整体升到熔点，最后全部为冰，
\[T_f=T_m+\frac{E+m_wL_f}{(m_i+m_w)c_i}<T_m.
\]
若 $E>m_iL_f$，冰全部熔化后还有多余热量，最后全部为水，
\[T_f=T_m+\frac{E-m_iL_f}{(m_i+m_w)c_w}>T_m.
\]
四个区间可用一个具体的思考过程串联起来。先把冰升到熔点、把水降到熔点，检查差额 $E$；有富余就用来熔冰，有不足就靠部分水凝固补足。只有当所有冰已经熔完仍有富余，才由全部液水共同升温；只有当所有水已经凝固仍有不足，才由全部固体一起降温。此时分母的总热容必须改用末态所属相，不能仍把初始冰、水的热容机械相加。

在 $E=0$ 处，没有净相变量；在 $E=m_iL_f$ 处，最后一份冰刚熔化；在 $E=-m_wL_f$ 处，最后一份水刚凝固。相邻分段在这些边界均给 $T_f=T_m$，可用来检查表达式的连续性。

这些结果由同一能量守恒分段得到。每算出一个末态，都要反查温度是否落在假定相态允许的范围内；如果把量热器热容计入，便需在同一温区的能量项中加入相应贡献。

\section{蒸发、饱和与有限水量}
\subsection{蒸发与凝结的动态平衡}
蒸发可在液体表面发生，表面分子中能量较大的部分逸出为气体；蒸气分子也会重新进入液体。封闭空间内，当单位时间内蒸发与凝结的物质量相等时，气液界面达到动态平衡，而不是分子停止交换。

纯液体与其蒸气在平直界面上平衡时，蒸气分压为该温度的饱和蒸气压 $p_s(T)$。在气液共存曲线范围内，$p_s$ 随温度升高通常迅速增大；给定温度下，液体多少主要决定能维持饱和状态到多大空间，并不直接改变平衡饱和压。

这一结论有明确条件：液体仍存在、相平衡已建立、界面曲率效应可以忽略。极小液滴的曲率、溶液成分以及显著非平衡蒸发，都可能使简单的 $p_v=p_s(T)$ 模型需要修正。


设封闭系统总计含 $n_w$ mol水，气体空间为 $V$，温度为 $T$，液体体积忽略，蒸气理想。在该空间保持饱和所需的蒸气量为

\[n_s(T,V)=\frac{p_s(T)V}{RT}.
\]
若 $n_w>n_s$，只有 $n_s$ mol在气相，其余仍为液体，蒸气分压为 $p_s$；若 $n_w<n_s$，全部水已为气体，分压为 $n_wRT/V<p_s$。两种情况可合写为

\begin{equation}\label{eq:t4vapour}
\boxed{n_v=\min\!\left(n_w,\frac{p_sV}{RT}\right)},\qquad
\boxed{p_v=\min\!\left(\frac{n_wRT}{V},p_s\right)}.
\end{equation}
等号处是液体刚尽或刚开始凝结的分界。这个判据适用于固定总水量、无其他气液反应且液体体积可忽略的模型；若容器内液体占据明显体积，须用实际气体空间 $V_g$ 代替总容积。

可把 $n_s$ 理解为气相在当前 $T,V$ 下的饱和容量。扩大空间时，这个容量按 $V$ 增加，液水蒸发补充气相；在液水尚存时，增加的是气相分子数而不是蒸气分压。液体耗尽后，再扩张就没有分子补充，分压才按 $1/V$ 下降。压缩时顺序反向：先按固定水量升压，达到饱和后多余分子凝结。这个叙述把“分压不变”与“气体量不变”区分开，二者在饱和分支上不能同时成立。

“水蒸气饱和”与“容器内一定看得到液水”也要区分。全部水恰好形成饱和蒸气时，液水可为零；继续膨胀便转入不饱和分支。例\ref{t4-2}的折点就是这一边界状态。

\fig{fig:h4:finitewater}{同温扩大气体空间：液水补充蒸气，液体刚尽后再无分子来源}{\begin{tikzpicture}

\foreach \j/\hh/\vv in {0/1.5/{V_*/2},1/3/{V_*},2/4.5/{3V_*/2}}{\begin{scope}[xshift={3.8*\j*1cm}]\draw[world](0,5)--(0,0)--(1.6,0)--(1.6,5);\fill[blue!4](0,0)rectangle(1.6,\hh);\draw[very thick](0,\hh)--(1.6,\hh);\node at(.8,5.5){$T$相同};\node at(.8,-.6){$V=\vv$};\end{scope}}
\fill[blue!35](.5,0)--(.8,.3)--(1.1,0)--cycle;\node[below,align=center]at(.8,-1){$n_v=n_l=n_w/2$\\$p_v=p_s$};\node[below,align=center]at(4.6,-1){$n_v=n_w,\ n_l=0$\\$p_v=p_s$};\node[below,align=center]at(8.4,-1){$n_v=n_w,\ n_l=0$\\$p_v=2p_s/3$};

\end{tikzpicture}}
图\ref{fig:h4:finitewater}中左图底部的小片液水只表示相态，体积忽略。气体空间由$V_*/2$增至$V_*$，蒸气量加倍而分压不变；中图恰好没有液体，也仍可达到饱和。继续扩至$3V_*/2$时总水量仍为$n_w$，分压才降为$2p_s/3$。同一套物质分配关系将对应下面的两条折线，不能只从压强平稳就认定气相量固定。



\fig{fig:phase-amounts}{定温改变气体空间时，气相与液相水量之和保持不变}{
\begin{tikzpicture}[x=3cm,y=2.2cm]
\draw[axis] (0,0)--(2.4,0) node[right]{$V/V_*$};
\draw[axis] (0,0)--(0,1.4) node[above]{物质的量$/n_w$};
\draw[world,blue] (.05,.05)--(1,1)--(2.2,1);
\draw[thick,dashed] (.05,.95)--(1,0)--(2.2,0);
\draw[guide] (1,0)--(1,1);\node[below] at (1,0){$1$};\node[left] at (0,1){$1$};
\node[above,blue] at (1.7,1){气相 $n_v/n_w$};
\node[above right] at (.15,.85){液相 $n_l/n_w$};
\node[below,align=center] at (.55,-.3){气液共存\\$p_v=p_s$};
\node[below,align=center] at (1.7,-.3){液体已尽\\$p_v<p_s$};
\end{tikzpicture}}
图\ref{fig:phase-amounts}与压强图描述同一个过程，横坐标 $V_*=n_wRT/p_s$ 由水量、温度共同决定。在共存段，$n_v/n_w=V/V_*$ 随体积线性增加，$n_l/n_w=1-V/V_*$ 同时减少；液体恰尽后，$n_v=n_w$、$n_l=0$，再膨胀不能继续增加气相水量。每一个状态都同时满足
\[
 n_v+n_l=n_w,\qquad n_l\ge0,\qquad p_v\le p_s(T).
\]
若计算得到负的剩余液量，说明错误地把共存分支延伸过了 $V_*$；若全部气态试算得到 $p_v>p_s$，则应考虑凝结分支。这两个检验分别排除两种相反的错误。含空气时图中的水量关系仍成立，空气只需另用自身守恒关系求分压。

以上分支图以温度已知为前提。若容器绝热或温度也待求，$p_s(T)$、气相量、液相量与温度须同时满足能量守恒，不能先把温度固定，再用一个与实际端点不相容的饱和压。决定相态、写物质守恒和写能量守恒应相互校验。

\subsection{含空气时，总压与蒸气分压分别由什么决定}
设不凝结空气量为 $n_a$，总水量为 $n_w$。理想混合气体满足
\[p=p_a+p_v,\qquad p_a=\frac{n_aRT}{V}.
\]
空气量固定，但水蒸气量可因相变改变。气液共存时 $p_v=p_s(T)$，因此总压 $p=n_aRT/V+p_s(T)$，一般不会因为“蒸气饱和”就保持不变。

定温膨胀时，液体刚尽的体积为 $V_*=n_wRT/p_s$。令 $u$ 表示当前气体体积，得到
\[\boxed{p(u)=\begin{cases}n_aRT/u+p_s,&u\le V_*,\\(n_a+n_w)RT/u,&u\ge V_* .\end{cases}}
\]
\fig{fig:sat}{空气与有限水量在定温下共同膨胀；总压和水蒸气分压在同一体积处切换分支}{\begin{tikzpicture}

\begin{scope}[x=2.2cm]
\draw[axis](0,0)--(2.9,0)node[right]{$V/V_*$};\draw[axis](0,0)--(0,3.5)node[above]{$p/p_s$};\draw[blue,thick](.25,1)--(1,1);\draw[blue,thick,domain=1:2.5,samples=90]plot(\x,{1/\x});\draw[red,thick,domain=.25:1,samples=90]plot(\x,{1+.4/\x});\draw[red,thick,domain=1:2.5,samples=90]plot(\x,{1.4/\x});\draw[guide](1,0)--(1,1.4);\fill(1,1.4)circle(.035);\fill[blue](1,1)circle(.035);\node[below]at(1,0){$1$};\node[left]at(0,1){$1$};\draw[red,thick](1.7,3.1)--(1.9,3.1)node[right]{总压};\draw[blue,thick](1.7,2.65)--(1.9,2.65)node[right]{蒸气分压};\node at(1.5,-1){$n_a/n_w=0.4$};\end{scope}

\end{tikzpicture}}

图\ref{fig:sat}的左分支中，蒸气分压恒定，但空气分压随膨胀下降；右分支中两组分的量均固定，总压按 $1/u$ 下降。两式在折点连续，斜率却改变。解题应同时检查压强连续和相态切换，不能把两段各算完却出现不合理的压力跳变。

对等温、组分理想且不考虑空气溶解的模型，水蒸气分压仍由式\eqref{eq:t4vapour}决定，空气分压为 $p_a=n_aRT/V$，总压为两者之和。只要还有液体，增大容积会使更多液体汽化，蒸气分压被相平衡固定；全部汽化后，水的气相物质的量不再增加，分压才按 $1/V$ 下降。


两段的切换位置由 $n_wRT/V_*=p_s(T)$ 决定，而不是由总压恰好等于饱和压决定。加入空气通常提高总压，却不把固定温度下的饱和蒸气压改成同一个更高数值。检查题目时应分别列水量守恒、空气量守恒和相平衡条件，最后才合成总压。


\section{湿度、冷却凝结与露点}
\subsection{相对湿度的定义及物理意义}
相对湿度定义为 $\varphi=p_v/p_s(T)$，通常用百分数表示。它比较实际水蒸气分压与该温度的饱和压；并不直接等于水蒸气在混合气体中的摩尔分数，也不等于蒸气质量占空气质量的比例。

例如相对湿度75\%表示实际蒸气压是当前饱和压的四分之三，不表示每100个气体分子中有75个水分子。水蒸气摩尔分数应是 $p_v/p$，分母是总压；相对湿度的分母则是 $p_s(T)$。两者量纲相同却描述不同关系，在含大量空气的容器中数值可以相差很大。

同一水蒸气分压在较高温度下可能相对湿度较低，因为该温度的饱和压较高。升温而暂不改变水量时，相对湿度往往下降；降温时饱和压下降较快，相对湿度增加，达到饱和后开始凝结。


刚性封闭容器冷却、尚未凝结时，$n_v,V$ 都不变，故 $p_v(T)=p_{vi}T/T_i$。开始凝结的温度满足

\[p_s(T_d)=p_{vi}\frac{T_d}{T_i}.
\]
若外界维持总压不变，且无凝结前各气体量比例不变，则水蒸气摩尔分数不变、分压也不变，此时露点近似满足 $p_s(T_d)=p_{vi}$。两个式子的区别来自过程条件，不能把定压冷却的关系原封不动用于刚性容器。

\fig{fig:dewpoint}{定容冷却时，试算蒸气分压与饱和压在凝结起点相交}{\begin{tikzpicture}[x=1cm,y=1cm]
\draw[axis] (0,0)--(5.4,0) node[right]{$T$};\draw[axis] (0,0)--(0,3.5) node[above]{$p_v$};\node[below left] at (0,0){$O$};\draw[guide] (0,0)--(5,2.5);\draw[world] plot coordinates {(.8,.12)(1.4,.28)(2,.55)(2.6,.9)(3,1.5)(3.4,2.2)(3.8,3.2)};\draw[ray] (3.6,1.8)--(3,1.5);\draw[ray] (3,1.5)--(2.6,.9)--(2,.55);\fill (3,1.5) circle(.05);\node[above left] at (3,1.5){$D$};\node[right] at (3.8,3.2){$p_s(T)$};\node[right] at (5,2.5){无凝结试算线};\node[below] at (3,0){$T_d$};
\end{tikzpicture}}

图\ref{fig:dewpoint}中从较高温度向左冷却，起初沿固定蒸气量的试算线；到 $D$ 点达到饱和，再冷却时部分蒸气凝结，实际分压转而沿 $p_s(T)$ 变化。凝结以后的蒸气量应重新用 $p_s(T)V/(RT)$ 求出，凝结量为初始总水量减去末态气相水量。

练习\ref{t4-10}采用的就是先试算、再切换的方法。若一开始就假定全部水蒸气量不变，会算出超过饱和压的末态；若一开始就强令饱和，也可能在水量不足时虚构凝结液。

\section{沸腾、液柱与多组分蒸气}
\subsection{蒸发与沸腾有何不同}
蒸发在表面发生，可在低于沸点的温度进行；沸腾则要求液体内部气泡能够持续存在并长大。气泡不只受蒸气压作用，还受外部液体压强、曲率张力和可能存在的空气分压影响。

深度 $h$ 处半径 $r$ 的气泡若近似球形，机械平衡可写为
\[p_s(T)+p_{a,\text{泡}}=p_{\text{表面}}+\rho gh+\frac{2\sigma}{r}.
\]
忽略泡内空气、液柱和表面张力后，才得到常用沸腾条件 $p_s(T)=p_{\text{外}}$。外压降低时，达到这一条件所需温度降低；液体深处或很小气泡需要克服额外压强，不能只比较表面外压。

表面水蒸气是否饱和与液体是否整体沸腾也不是同一条件。一个敞口水面附近可以接近饱和，液体内部却因外压仍较大而不能形成持续长大的纯蒸气泡。沸点的判据针对泡内外的机械平衡；饱和压的定义针对气液之间的相平衡。将两条条件在同一气泡上联合，才得到沸腾温度，而不能只见“有饱和蒸气”就判断正在沸腾。

若泡中有空气，气泡膨胀时空气分压随体积增大而减小；与液体同温且充分交换的水蒸气分压则由温度决定，并不必然随膨胀降低。把空气和蒸气混为同一种固定气体，是含气泡问题的常见错误。


\fig{fv:h4:boiling}{体内气泡能否增长，取决于蒸气压能否抵住局部外压}{\begin{tikzpicture}[scale=.9]
\draw[thick](0,3.1)--(0,0)--(4.3,0)--(4.3,3.1);\fill[blue!9](0,0)rectangle(4.3,2.5);\draw[blue](0,2.5)--(4.3,2.5);
\draw[thick](2.1,1)circle(.5);\node at(2.1,1.2){$p_{\rm in}$};\draw[dim](2.1,1)--(2.6,1)node[midway,below]{$r$};\draw[->](1.1,1)--(1.6,1);\draw[->](3.1,1)--(2.6,1);
\draw[<->](4.8,1)--(4.8,2.5)node[midway,right]{$h$};\node at(2.2,3.5){$p_0$};\end{tikzpicture}}
图\ref{fv:h4:boiling}中，气泡中心深度$h$决定周围液体压强，半径$r$决定曲率附加压强。较小的泡需要更大的内外压差，不能只看液面处的大气压。若泡内含不凝结空气，同温扩张会降低其空气分压；水蒸气分压则在相平衡条件下仍由$T$决定。两种组分对体积变化的响应不同，必须分别处理。


\subsection{饱和蒸气与液柱平衡须同时满足}
倒扣管、封气管中，液柱关系使用总压，而空气守恒关系只用于空气分压。正确顺序是先画液面高度，写总压与外压的静水差，再从总压减去蒸气分压，最后对不凝结气体使用状态方程。

加热过程中，空气膨胀与饱和压增加会共同推动液面。有些装置的液面会由槽面以上降到槽面以下；方向不应预先假定，应由空气分压必须为正及气液平衡方程判定。例\ref{t4-3}还需检查管口浸没深度是否足够，防止空气逸出使原守恒关系失效。

恒温水银封气管中，水银蒸气分压恒定，但空气分压随气柱长度变化。竖管的开口方向决定水银柱压强加号或减号。例\ref{t4-4}通过明确“开口向下”，把几何方向与压强关系对应起来。

\subsection{含空气的压力阈值与互不相溶液体}
加热含空气、水和蒸气的封闭容器时，总压为 $p_a(T)+p_s(T)$。刚性且空气量不变时，空气分压正比于绝对温度。因而同一总压阈值可以在低于纯蒸气对应温度时达到，不能把压力阈值直接当作液体温度计。

若随后排出气体并维持总压，空气分压下降，水蒸气分压可以随着温度增加而增加；在液水仍足量时，温度逐渐接近纯蒸气在该总压下的饱和温度。例\ref{t4-5}用分压曲线区分首次达到阈值和后续状态。

两种互不相溶液体都能向同一气泡供给蒸气时，气泡总蒸气压约为 $p_{s1}(T)+p_{s2}(T)$。忽略液柱和曲率压差，持续生长条件为

\[p_{s1}(T)+p_{s2}(T)=p_0.
\]
因为每个分压都小于总压，该温度可低于两种纯液体各自的正常沸点。界面附近同时得到两种蒸气，可能比仅接触一种液体的区域更容易达到条件；实际最先沸腾位置还受成核、深度及供质速率影响。


气液共存曲线在临界点终止，临界温度以上不能单靠压缩形成明确的气液分离。此时也不能继续把“压强达到某个饱和值便凝结”的低温图像延伸使用。

即便温度低于临界温度，理想气体方程在高密度、接近凝结区域也只是近似。本讲采用它来处理稀薄蒸气和分压关系；题目若给出实际状态数据，应按数据与适用模型判断，而不是要求真实蒸气严格服从 $pV=nRT$。

\section{不同温度容器之间的蒸气交换}
\subsection{机械平衡与热平衡不应混同}
两个容器连通后可有相同压强，而温度因各自恒温装置作用保持不同。这是受外界维持的条件，不能再要求两边温度相等，也不能对两容器气体直接列无热交换的内能守恒。

忽略管容积、静压差及显著流动后，若共同压强为 $p$，气相总物质的量为
\[n_g=\frac pR\left(\frac{V_A}{T_A}+\frac{V_B}{T_B}\right).
\]
同压不代表两室气体量相等，比例为 $(V_A/T_A):(V_B/T_B)$。这一关系同时用于第一次凝结后的气相分配和第二次连通后的试算。


对纯水蒸气，两端若都要求饱和而温度不同，所需饱和压一般不同，与共同压强矛盾。若水量足够，冷端先达到自身饱和压并凝结；热端在同一较低压强下通常不饱和。

但“冷端决定压强”仍有水量条件。应先假设全部气态，由总水量和上一公式求共同压强；若超过冷端饱和压，才切换到冷端凝结分支，差额物质的量凝结为液体。

\fig{fig:cold-trap}{冷热容器的共同压强及冷端凝结条件}{\begin{tikzpicture}[x=1cm,y=1cm]
\draw[world] (0,0) rectangle (3,2) (5,0) rectangle (8,2);\draw[world] (3,1)--(5,1);\draw[world](3.8,.85)--(4.2,1.15)--(4.2,.85)--(3.8,1.15)--cycle;\node[align=center] at (1.5,1){热端 $T_A$\\$p<p_s(T_A)$};\fill[blue!15] (5.02,.02) rectangle (7.98,.35);\node[align=center] at (6.5,1.2){冷端 $T_B<T_A$\\$p=p_s(T_B)$};\node at (4,-.7){水量足够时：冷端有凝结液，热端不饱和};
\end{tikzpicture}}

再次操作前，应按图\ref{fig:cold-trap}把水分为热端蒸气、冷端蒸气和凝结液，逐项确定实际保留量。关闭阀门只中断交换，不会使热端恢复为最初的饱和状态。若只抽走冷端气体，剩余凝结液还会蒸发补充；若把气体和凝结液都移除，第二次连通的总水量就更少。例\ref{t4-6}须以操作后的保留水量重新试算饱和条件，而不能沿用第一次的总量或压强。

\section{小扰动稳定性与绝热稳流}
\subsection{气体的恢复作用取决于扰动时间尺度}
把气体从平衡体积$V$微小改变为$V+\delta V$。等温关系$pV=\text{常量}$的微分为$V\dd p+p\dd V=0$，于是$\dd p/\dd V=-p/V$，一阶变化为$\delta p\simeq-p\delta V/V$。可逆绝热关系$pV^\gamma=\text{常量}$的微分为$V^\gamma\dd p+p\gamma V^{\gamma-1}\dd V=0$，约去$V^\gamma$后得到$\delta p\simeq-\gamma p\delta V/V$。两者的一阶响应可统一写成$\delta p=-\kappa p\delta V/V$，其中慢扰动充分换热时常取$\kappa=1$，快扰动近似绝热时取$\kappa=\gamma$。

使用哪个响应系数由扰动时间与热交换时间比较决定，不能只看缸壁材料名称。对实际存在耗散或热滞后的系统，这两种近似之间还可能有更复杂的动力学行为。


管口液柱恰到边缘时，向上位移 $x$ 使气柱增加 $Sx$，液体溢出后负载液柱降低 $x$。因此
\[\delta p_g=-\frac{\kappa p}{\ell}x,\qquad\delta p_l=-\rho gx,
\quad\delta F=S\left(\rho g-\frac{\kappa p}{\ell}\right)x.
\]
\fig{fig:h4:overflowresponse}{只考察向上溢液的微小位移；气压下降与负载下降的竞争决定单侧恢复性}{\begin{tikzpicture}

\draw[world](0,4.2)--(0,0)--(1,0)--(1,4.2);\fill[blue!6](0,0)rectangle(1,2.5);\fill[gray!25](0,2.5)rectangle(1,4);\draw[guide](-.2,2.5)--(1.5,2.5);\draw[blue,thick](0,2.8)--(1,2.8);\draw[flow](.5,2.5)--(.5,2.8);\node[right]at(1.2,2.65){$x>0$};\draw[flow](1,4)--(1.8,4)node[right]{溢出};\draw[dim](-.6,0)--(-.6,2.5)node[midway,left]{$\ell$};
\begin{scope}[xshift=5cm,yshift=2cm]\draw[axis](0,0)--(3.8,0)node[right]{$x$};\draw[axis](0,-1.5)--(0,1.8)node[above]{$\delta F$};\draw[red,thick](0,0)--(3,1.2);\draw[blue,thick](0,0)--(3,-1.2);\node[red,above]at(2,1.3){$\kappa p/\ell<\rho g$};\node[blue,below]at(2,-1.3){$\kappa p/\ell>\rho g$};\end{scope}

\end{tikzpicture}}
图\ref{fig:h4:overflowresponse}左图的蓝线是气液界面上移后的示意位置，管口高度不变，因而上方液柱变短。右图只画$x>0$分支：若气体压降更大，净力指向下方；若负载压降更大，净力继续推向上方。这是溢出后的单侧判据，不能把已流失的液体当作仍能在负位移时自动补回。



若 $\kappa p/\ell>\rho g$，上移后净力向下，具有恢复趋势；若小于则上移被进一步放大。等号时一阶项消失，需看高阶项与实际边界。向下扰动可能不溢液，负载变化规律不同，须另列分支，不能把上移关系直接延伸。

对缓慢等温溢出，平衡温度可写为$T(h)=C(p_0+\rho gh)(L-h)$。对两个因子分别求导，
\[
T'(h)=C\bigl[\rho g(L-h)-(p_0+\rho gh)\bigr].
\]
令其为零，得$\rho g(L-h)=p_0+\rho gh$。此时$\ell=L-h$、$p=p_0+\rho gh$，所以极值条件正是$p/\ell=\rho g$，即恢复系数变为零；而$T''(h)=-2C\rho g<0$确认温度取极大值。这把第二讲的临界温度方法与本讲的小位移方法联系起来。

局部存在恢复力不等于实际系统会无耗散地周期振动。已经溢出的液体不会自动回到管中，物质量和后续平衡位置可能改变。例\ref{t4-8}及练习\ref{t4-13}特别区分这两个结论。

\subsection{喷管能量式中的焓与出口速度}
稳态意味着固定位置的宏观量不随时间变化，并不表示流体质点的温度和速度不变。气体沿喷管运动时，压力能和内能可以转化为宏观动能，入口、出口状态由流动关系联系。

对稳态绝热、无轴功、势能变化可忽略的喷管，单位质量能量守恒为
\[\frac{C_pT_1}{\mu}+\frac{v_1^2}{2}=\frac{C_pT_2}{\mu}+\frac{v_2^2}{2}.
\]
这里 $C_p$为摩尔热容、$\mu$为kg/mol，$C_pT/\mu$才是单位质量焓。第三讲已将内能与推动功合并为焓；在喷管中，上游流体对气团做入口压力功，气团对下游做出口压力功，两端都须计入。因此方程使用入口焓减出口焓，不能只取 $C_VT/\mu$，也不能在使用焓后再加一次压力功。刚性管壁不做位移功，并不妨碍流体之间通过压强作用交换能量。

若流动还可逆、气体量随质点守恒且热容恒定，等熵关系给
\[\frac{T_2}{T_1}=\left(\frac{p_2}{p_1}\right)^{R/C_p}.
\]
入口速度可忽略时，联立得
\[\boxed{v_2=\sqrt{\frac{2C_pT_1}{\mu}\left[1-\left(\frac{p_2}{p_1}\right)^{R/C_p}\right]}}.
\]
\fig{fig:h4:nozzle}{喷管控制面两端都有压力功；焓的减少可以转化为流体宏观动能}{\begin{tikzpicture}

\draw[world](0,2.2)..controls(1.2,2.2)and(2,1.45)..(2.7,1.45)..controls(3.7,1.45)and(4.7,2.1)..(6,2.1);\draw[world](0,0)..controls(1.2,0)and(2,.75)..(2.7,.75)..controls(3.7,.75)and(4.7,.1)..(6,.1);\draw[guide,rounded corners](-.25,-.25)rectangle(6.25,2.5);
\draw[flow](-.9,1.25)--(.7,1.25)node[midway,above]{$v_1$};\draw[flow](5.4,.7)--(7.2,.7)node[midway,below]{$v_2$};\draw[force](-1,.4)--(0,.4)node[midway,below]{$p_1S_1$};\draw[force](7.1,1.6)--(6,1.6)node[midway,above]{$p_2S_2$};\node at(.3,3.1){$p_1,T_1$};\node at(5.8,3.1){$p_2,T_2$};\node at(3,-1.3){$C_pT_1/\mu+v_1^2/2=C_pT_2/\mu+v_2^2/2$};

\end{tikzpicture}}
图\ref{fig:h4:nozzle}的压力箭头都表示外侧流体对控制面内流体的作用：入口推动流体进入，出口压力阻碍流体流出。两者分别与随流体携带的内能合成入口、出口焓；图中管形只示意流路，具体面积比和马赫数仍需另由流动条件确定。使用焓之后，不能再把同一入口推动功多加一次。



应核对三点：出口压强降低使温度降低、速度增加；$p_2\to p_1$ 时速度趋零；根号内量纲为 $\mathrm{m^2/s^2}$。例\ref{t4-9}把出口压强作为可实现的等熵状态给定，具体喷管还可能涉及喉部阻塞、面积比和激波，不能仅由这个能量式确定质量流率或所有出口条件。

